% =========================================================================
% IAIFI Summer Workshop 2026 -- 12 minute slot
% Re-cut of talk_iaifi_2026.tex:
%   * result moved to slide 2
%   * old slides 5+6+7 merged; old 8+9 merged; RG table moved to backup
%   * overlays added; no \animategraphics (works in every PDF viewer)
%   * 8pt typographic floor; phase colors no longer overloaded
% =========================================================================
\documentclass[aspectratio=169,10pt]{beamer}

\usetheme{default}
\useinnertheme{rectangles}
\setbeamertemplate{navigation symbols}{}
\setbeameroption{hide notes}
\setbeamersize{text margin left=0.7cm, text margin right=0.7cm}

\usepackage{amsmath,amssymb}
\usepackage[T1]{fontenc}
\usepackage[scaled=0.96]{FiraSans}
\usepackage{newtxsf}
\renewcommand{\familydefault}{\sfdefault}
\usepackage{microtype}
\usepackage{booktabs}
\usepackage{graphicx}
\usepackage{tikz}
\usetikzlibrary{arrows.meta,positioning,calc}
\usepackage{pgfplots}
\usepgfplotslibrary{groupplots}
\pgfplotsset{compat=1.18}

% -------------------------------------------------------------------------
% ANIMATION TOGGLE.  Build both with scripts/build_talk.sh:
%   _animated.pdf : \animategraphics, autoplay. Needs Acrobat / pdf-presenter.
%   _static.pdf   : the same frames as click-through stills. Works everywhere
%                   (Preview, Skim, browsers, screen share).
% Present from _animated if you control the machine and it has Acrobat;
% carry _static as the fallback.
% -------------------------------------------------------------------------
\providecommand{\useanimations}{0}
\newif\ifanim
\ifnum\useanimations=1 \animtrue\else\animfalse\fi
\ifanim\usepackage{animate}\fi

% Overlay step at which slide 4's conclusion box appears: right after the
% animation in animated mode, after the fourth still in static mode.
\ifanim
  \newcommand{\manifoldbox}{2}
\else
  \newcommand{\manifoldbox}{4}
\fi

% Slide measures are short: avoid mid-word hyphenation.
\hyphenpenalty=10000
\exhyphenpenalty=10000

% -------------------------------------------------------------------------
% CONFIRM BEFORE PRESENTING.
% The old deck's speaker note cited 25 runs / 75 epochs for the primary MLP
% scheduling experiment, but main.tex Table (high-overfitting regime) lists
% 20 runs / 50 epochs. These are different experiments or a stale note.
% Set these to whatever the h-sweep run actually used.
% -------------------------------------------------------------------------
\newcommand{\mlpruns}{20}
\newcommand{\epochsfast}{29--31}
\newcommand{\epochsbase}{75}

\definecolor{ink}{HTML}{202121}
\definecolor{forest}{HTML}{0E3B2C}
\definecolor{gold}{HTML}{D8AD62}
\definecolor{gilt}{HTML}{9E7B38}
\definecolor{paper}{HTML}{F4EFE3}
\definecolor{mist}{HTML}{C9CBC7}
\definecolor{muted}{HTML}{60645F}
% Phase colors. Reserved meanings, used nowhere else in the deck:
%   orderedblue = ordered phase, forest = critical edge, clay = chaotic phase.
\definecolor{clay}{HTML}{B04A2F}
\definecolor{orderedblue}{HTML}{3A6EA5}
% Caveats and scope limits get their own accent so clay stays "chaotic".
\definecolor{caveat}{HTML}{7A5C2E}

\setbeamercolor{normal text}{fg=ink,bg=paper}
\setbeamercolor{background canvas}{bg=paper}
\setbeamercolor{frametitle}{fg=forest}
\setbeamercolor{structure}{fg=forest}
\setbeamercolor{block title}{fg=white,bg=forest}
\setbeamercolor{block body}{fg=ink,bg=paper!85!mist}
\setbeamercolor{itemize item}{fg=gilt}
\setbeamertemplate{itemize item}{\raisebox{0.15em}{\tiny$\blacktriangleright$}}
\setlength{\abovedisplayskip}{0.4em}
\setlength{\belowdisplayskip}{0.4em}

\newcommand{\relu}{\mathrm{ReLU}}
\newcommand{\E}{\mathbb{E}}
\newcommand{\cmulogo}{assets/logos/cmu_seal.png}
\newcommand{\iaifilogo}{assets/logos/IAIFI-logo-black.pdf}
\newcommand{\regularbackground}{assets/iaifi_regular_background.png}

% Caveat text: one consistent style, never below 8pt.
\newcommand{\scope}[1]{{\footnotesize\color{caveat}#1}}
\newcommand{\srcnote}[1]{{\footnotesize\color{muted}#1}}

% -------------------------------------------------------------------------
% Photos for the broken-telephone slide. Drop these three files in and they
% appear automatically; until then the slide falls back to labelled boxes,
% so the deck always compiles.
%   assets/cat.png     the cat photo
%   assets/dog.png     the dog photo
%   assets/static.png  TV static, used for BOTH outputs (they end up
%                      indistinguishable, which is the whole point)
% Any raster format pdflatex accepts works; rename to .png or edit below.
% -------------------------------------------------------------------------
\newcommand{\catimg}{assets/cat.png}
\newcommand{\dogimg}{assets/dog.png}
\newcommand{\staticimg}{assets/static.png}
% #1 file, #2 height, #3 frame colour, #4 fallback label
\newcommand{\telephoto}[4]{%
  \IfFileExists{#1}%
    {\includegraphics[height=#2]{#1}}%
    {\parbox[c][#2][c]{1.25cm}{\centering\scriptsize\color{#3}#4}}%
}
% One input under a particular dropout mask: the photo, some static, and a
% set of blanked blocks standing in for the dropped units. Coordinates in
% #4 are fractions of the image, so the same list works at any size.
% #1 file, #2 size, #3 noise opacity, #4 list of x/y/w/h blocks
\newcommand{\maskedphoto}[4]{%
  \begin{tikzpicture}[inner sep=0pt,outer sep=0pt,x=#2,y=#2]
    \node[inner sep=0pt,anchor=south west] at (0,0)
      {\includegraphics[width=#2,height=#2]{#1}};
    \node[inner sep=0pt,anchor=south west,opacity=#3] at (0,0)
      {\includegraphics[width=#2,height=#2]{\staticimg}};
    \foreach \bx/\by/\bw/\bh in {#4} {
      \fill[paper!92!ink] (\bx,\by) rectangle ({\bx+\bw},{\by+\bh});
    }
  \end{tikzpicture}%
}

% Square-cropped variant: static.png is 4:3, so force it to the same
% footprint as the square photos.
\newcommand{\squarephoto}[3]{%
  \IfFileExists{#1}%
    {\includegraphics[width=#2,height=#2]{#1}}%
    {\parbox[c][#2][c]{#2}{\centering\scriptsize\color{muted}#3}}%
}
% Partially-degraded image: the photo with static laid over it at opacity
% #3. Gives the mid-depth state without needing extra image files.
% #1 file, #2 height, #3 noise opacity
\newcommand{\noisyphoto}[3]{%
  \IfFileExists{#1}%
    {\begin{tikzpicture}[inner sep=0pt,outer sep=0pt]
       \node[inner sep=0pt] (np) {\includegraphics[height=#2]{#1}};
       \node[inner sep=0pt,opacity=#3] at (np.center)
         {\includegraphics[height=#2,width=#2]{\staticimg}};
     \end{tikzpicture}}%
    {\parbox[c][#2][c]{#2}{\centering\scriptsize\color{muted}noisy}}%
}

% The prebuilt background PNG has an OLD CMU seal and the gold rule baked in,
% which double-printed under the foreground chrome below. Just paint the paper
% colour; the hook draws the rule and the (current) seal.
\setbeamertemplate{background canvas}{%
  \color{paper}\rule{\paperwidth}{\paperheight}%
}
% Slide chrome in the shipout foreground. With \animategraphics gone there are
% no optional-content groups left to swallow it, so the old per-frame fallback
% redraws are no longer needed.
\makeatletter
\AddToHook{shipout/foreground}{%
  \ifbeamer@plainframe\else
    \begin{tikzpicture}[remember picture,overlay]
      \node[anchor=north east]
        at ($(current page.north east)+(-0.62cm,-0.15cm)$)
        {\includegraphics[height=1.02cm]{\cmulogo}};
      \draw[gold,line width=1.2pt]
        ($(current page.north west)+(0.70cm,-1.14cm)$) --
        ($(current page.north west)+(3.30cm,-1.14cm)$);
      \node[anchor=south west,text=muted,font=\scriptsize]
        at ($(current page.south west)+(0.70cm,0.10cm)$)
        {Lucas Fernandez-Sarmiento};
      \node[anchor=south,text=muted,font=\scriptsize]
        at ($(current page.south)+(0,0.10cm)$)
        {IAIFI Workshop 2026};
      \node[anchor=south east,text=muted,font=\scriptsize]
        at ($(current page.south east)+(-0.70cm,0.10cm)$)
        {\insertframenumber};
    \end{tikzpicture}%
  \fi
}
\makeatother

\setbeamertemplate{background}{}

\setbeamertemplate{frametitle}{%
  \nointerlineskip
  \vspace{0.20cm}
  \hbox to \textwidth{%
    \vbox to 0.60cm{\vfil
      \hbox{{\color{forest}\bfseries\Large\insertframetitle}}%
      \vfil}%
    \hfill%
  }%
  \vspace{0.03cm}
  \hbox to \textwidth{%
    \rule{0pt}{1.2pt}\hfill%
  }%
  \vspace{0.08cm}
}

\setbeamertemplate{footline}{%
  \leavevmode
  \hbox to \paperwidth{\rule[-0.11cm]{0pt}{0.36cm}\hfill}%
}

\title[Dropout Universality]{Dropout Universality}
\subtitle{Critical Scaling and Optimal Scheduling at the Edge-of-Chaos}
\author{Lucas Fernandez-Sarmiento}
\institute{Carnegie Mellon University}
\date{August 12, 2026}

\begin{document}

% =========================================================================
% 1. Title  (~20 s)
% =========================================================================
\begin{frame}[plain]
  \begin{tikzpicture}[remember picture,overlay]
    \fill[paper] (current page.south west) rectangle (current page.north east);
    \draw[gold,line width=0.8pt]
      ($(current page.north west)+(0.75,-2.22)$) --
      ($(current page.north east)+(-0.75,-2.22)$);
    \node[anchor=north west] at ($(current page.north west)+(0.75,-0.38)$)
      {\includegraphics[height=1.02cm]{\iaifilogo}};
    \node[anchor=west,text=forest] at ($(current page.north west)+(1.88,-0.87)$)
      {\bfseries 2026 IAIFI Summer Workshop};
    \node[anchor=north east] at ($(current page.north east)+(-0.62,-0.18)$)
      {\includegraphics[height=1.92cm]{\cmulogo}};
  \end{tikzpicture}

  \vspace{2.48cm}
  \begin{center}
    {\color{ink}\bfseries\fontsize{22}{27}\selectfont
    Dropout Universality}\\[0.35em]
    {\color{forest}\fontsize{15}{19}\selectfont
    Critical Scaling and Optimal Scheduling at the Edge-of-Chaos}

    \vspace{0.65cm}
    {\color{gold}\rule{6.0cm}{1.0pt}}

    \vspace{0.48cm}
    {\Large\color{ink} Lucas Fernandez-Sarmiento}\\[0.18em]
    {\large\color{muted} Carnegie Mellon University}

    \vspace{0.28cm}
    {\small\color{forest}\bfseries ICML 2026}
    \ {\small\color{muted}\(\cdot\)}\
    {\small\color{muted}\href{https://arxiv.org/abs/2605.21648}{arXiv:2605.21648}}

    \vfill
    {\color{muted}August 12, 2026 \(\cdot\) IAIFI \(\cdot\) MIT}
    \vspace{0.4cm}
  \end{center}
  \note{Say the one-liner here: dropout is an external field, and at fixed
  budget you should put it early. 20 seconds, then move.}
\end{frame}

% =========================================================================
% 2. The standard intuition, and what it cannot tell you  (~50 s)
% =========================================================================
\begin{frame}{A usual suspect of regularization: dropout}
  \begingroup
  \small
  \begin{center}
  \makebox[\textwidth][c]{%
  \begin{tikzpicture}[x=1cm,y=1cm,>=Latex,scale=1.16,
                      every node/.style={transform shape}]
    % one input, masked afresh at every layer: the pattern changes and the
    % damage compounds with depth.  Panels are 30% larger than the original
    % pass; the gaps were squeezed to keep the total width inside the frame.
    \draw[orderedblue,line width=1.3pt,-{Latex[length=2.0mm]}]
      (1.58,0)--(3.68,0);
    \draw[orderedblue,line width=1.3pt,-{Latex[length=2.0mm]}]
      (5.18,0)--(7.28,0);
    \draw[orderedblue!55,line width=1.3pt,-{Latex[length=2.0mm]}]
      (8.78,0)--(9.58,0);
    \node[font=\small,text=muted] at (9.95,0) {\(\cdots\)};

    \foreach \x in {2.22,2.92,5.82,6.52} {
      \node[draw=forest!60,fill=forest!8,minimum width=0.30cm,
            minimum height=2.35cm,rounded corners=1pt] at (\x,0) {};
    }

    \node[draw=orderedblue,line width=0.9pt,rounded corners=2pt,
          inner sep=1pt,fill=paper] at (0.80,0)
          {\telephoto{\catimg}{1.55cm}{orderedblue}{cat}};

    \node[draw=orderedblue!60,line width=0.9pt,rounded corners=2pt,
          inner sep=1pt,fill=paper] at (4.40,0)
      {\maskedphoto{\catimg}{1.55cm}{0.26}%
        {0.08/0.60/0.24/0.19, 0.44/0.12/0.19/0.24,
         0.70/0.64/0.21/0.22, 0.31/0.36/0.13/0.13}};

    \node[draw=orderedblue!40,line width=0.9pt,rounded corners=2pt,
          inner sep=1pt,fill=paper] at (8.00,0)
      {\maskedphoto{\catimg}{1.55cm}{0.52}%
        {0.55/0.55/0.27/0.23, 0.06/0.13/0.23/0.21,
         0.33/0.74/0.20/0.18, 0.76/0.22/0.17/0.17,
         0.14/0.42/0.15/0.16, 0.60/0.02/0.18/0.14,
         0.40/0.40/0.14/0.15}};

    \foreach \x/\lab in {0.80/{input}, 4.40/{a few layers in},
                         8.00/{deeper still}} {
      \node[font=\scriptsize\color{muted}] at (\x,-1.40) {\lab};
    }
  \end{tikzpicture}}
  \end{center}

  \vspace{0.05em}
  \begin{center}
    {\footnotesize\color{muted}A fresh mask at every layer: the pattern
    changes and the damage compounds.}
  \end{center}

  \vspace{0.20em}
  \begin{columns}[T,onlytextwidth]
    \begin{column}{0.48\textwidth}
      {\color{forest}\bfseries The standard theory}

      \vspace{0.10em}
      Usual NNs are overparametrized, so there are enough parameters to memorize training neural patterns ("coadaptation"). Dropout blocks activations randomly, mitigating this. 

      \vspace{0.15em}
      \srcnote{Srivastava et al., JMLR \textbf{15}, 1929 (2014).}
    \end{column}

    \begin{column}{0.48\textwidth}
      \onslide<2->{%
      {\color{clay}\bfseries The practice}

      \vspace{0.10em}
      People usually do a grid search over the probability of dropping out and call it a day.

      \vspace{0.40em}
      \begin{tikzpicture}[scale=0.85,every node/.style={transform shape}]
        \node[
          fill=forest!10!paper, fill opacity=0.45,
          draw=forest, line width=0.9pt,
          text opacity=1, text=ink,
          rounded corners=5pt, font=\normalsize,
          inner xsep=0.30cm, inner ysep=0.22cm,
          text width=0.74\linewidth, align=left
        ] {Reasonable. But is it \textbf{optimal}? Can we do better?};
      \end{tikzpicture}}
    \end{column}
  \end{columns}
  \endgroup
  \note{50 s. Give the textbook reading its due in one breath, then say
  plainly what it cannot deliver. The last line hands off to the
  collective-variable slide.}
\end{frame}

% =========================================================================
% 3. THE RESULT, UP FRONT  (~60 s)
% =========================================================================
% Built in eight beats: the headline alone, then it retreats into the
% frame title; the constant schedule at full size, an arrow, the step
% schedule it uncovers; the pair settles into the left column; the
% numbers get a beat at full size before landing on the right.
%
% \schedfig{scale}{show arrow}{show step} -- one geometry, used big for the
% centre-stage beats and small for the final layout, so nothing is redrawn
% differently between the two.
\newcommand{\schedfig}[3]{%
  \begin{tikzpicture}[x=1cm,y=1cm,scale=#1,
                      every node/.style={transform shape}]
    % frozen box: the step row must not shove the constant row when it lands
    \useasboundingbox (-0.10,-0.95) rectangle (5.30,2.30);
    % --- constant schedule ---
    \node[anchor=west,font=\footnotesize\bfseries,text=gilt]
      at (-0.05,1.95) {constant};
    \foreach \i in {0,...,11}{
      \fill[gilt!70] ({0.42*\i},1.25) rectangle ({0.42*\i+0.30},1.67);
    }
    \ifnum#2=1
      \draw[clay,line width=1.5pt,-{Latex[length=2.6mm]}]
        (2.46,1.16) -- (2.46,0.76);
    \fi
    % --- front-loaded: half the layers (6 of 12), twice the height.
    %     Equal area, so "same total dropout" is literally true. This is
    %     the paper's Step schedule: f = 1/2, h_max = 0.2, hbar = 0.1.
    \ifnum#3=1
      \node[anchor=west,font=\footnotesize\bfseries,text=forest]
        at (-0.05,0.95) {front-loaded};
      \foreach \i in {0,...,5}{
        \fill[forest] ({0.42*\i},-0.20) rectangle ({0.42*\i+0.30},0.64);
      }
      \foreach \i in {6,...,11}{
        \draw[forest!30,line width=0.4pt]
          ({0.42*\i},-0.20) rectangle ({0.42*\i+0.30},-0.18);
      }
    \fi
    \draw[muted,line width=0.6pt,-{Latex[length=1.6mm]}]
      (-0.02,-0.42) -- (5.20,-0.42);
    \node[anchor=north,font=\footnotesize,text=muted] at (2.55,-0.50)
      {layer \(\ell\)};
  \end{tikzpicture}%
}

% The numbers, written once: shown scaled up for their own beat, then at
% size in the right column.
\newcommand{\tldrmetrics}{%
  {\color{forest}\bfseries Final-epoch test loss}\\[0.10em]
  {\fontsize{17}{20}\selectfont\color{clay}\bfseries 18--35\%} lower
  {\small(MLP)}\qquad
  {\fontsize{17}{20}\selectfont\color{clay}\bfseries 4--6\%} lower
  {\small(ViT)}

  \vspace{0.40em}
  {\color{gold}\rule{6.8cm}{1.0pt}}

  \vspace{0.40em}
  {\color{forest}\bfseries Epochs to reach constant-dropout accuracy}\\[0.10em]
  {\fontsize{17}{20}\selectfont\color{clay}\bfseries \epochsfast}
  \ vs \ \epochsbase
  \ \ {\small\color{clay}\bfseries(\(\approx\)60\% fewer)}%
}

\begin{frame}
  \frametitle{\only<2->{TLDR: Front-load your dropout!}}
  \begingroup
  \small

  % ---- beat 1: the headline, alone -------------------------------------
  \only<1>{%
    \vfill
    \begin{center}
      {\color{forest}\bfseries\fontsize{34}{41}\selectfont
      Front-load your dropout!}
    \end{center}
    \vfill}

  % ---- beats 2-4: the schedules, centre stage ---------------------------
  \only<2>{\vfill\begin{center}\schedfig{1.55}{0}{0}\end{center}\vfill}
  \only<3>{\vfill\begin{center}\schedfig{1.55}{1}{0}\end{center}\vfill}
  \only<4>{\vfill\begin{center}\schedfig{1.55}{1}{1}\end{center}\vfill}

  % ---- beats 5-8: the settled layout ------------------------------------
  \only<5->{%
  \begin{columns}[c,onlytextwidth]
    \begin{column}{0.44\textwidth}
      \centering
      \schedfig{1.06}{0}{1}

      \vspace{0.55em}
      \onslide<8->{%
        {\footnotesize\color{muted}Same total dropout. Same cost per epoch.}}
    \end{column}

    \begin{column}{0.52\textwidth}
      % \onslide (not \only) so the column reserves its height from beat 5
      % and nothing below it shifts when the numbers arrive.
      \onslide<7->{\tldrmetrics}
    \end{column}
  \end{columns}

  \vspace{0.45em}
  \onslide<8->{%
  \begin{center}
    \begin{tikzpicture}
      \node[fill=forest!10!paper, fill opacity=0.45,
            draw=forest, line width=0.9pt,
            text opacity=1, text=ink,
            rounded corners=5pt,
            inner xsep=0.42cm,inner ysep=0.22cm]
        {Follows from applying principle of tuning to
         \textbf{criticality} to dropout distribution.};
    \end{tikzpicture}
  \end{center}}}

  % ---- beat 6: the numbers get their own moment, floated over the layout
  % so the settled columns underneath do not move when it goes away.
  \only<6>{%
    \begin{tikzpicture}[remember picture,overlay]
      % paper fill so the blow-up never collides with the schematic that
      % has already settled into the left column
      \node[anchor=center,inner sep=0.28cm,fill=paper,rounded corners=4pt]
        at ($(current page.center)+(3.30cm,0.15cm)$)
        {\scalebox{1.12}{\parbox{0.52\textwidth}{\small\tldrmetrics}}};
    \end{tikzpicture}}
  \endgroup
  \note{60 s. Eight clicks, so keep moving: headline, constant, arrow,
  step, land it, numbers, land them, the caveat and the promise. Land the
  number, promise the derivation, move on. Do not defend the experiment
  here; the results slide does that.}
\end{frame}

% =========================================================================
% 4. In search of good variables  (~40 s)
% =========================================================================
\begin{frame}{In search of better variables}
  \begingroup
  \small
  \begin{columns}[c,onlytextwidth]
    \begin{column}{0.600\textwidth}
      \vspace{0.35em}
      {\color{ink}\itshape\fontsize{14}{17.5}\selectfont
      ``With four parameters I can fit an elephant, and with five I can
      make him wiggle his trunk.''}\\[0.35em]
      {\color{gilt}\normalsize von Neumann to Fermi, via Dyson}

      \vspace{1.6em}
      \begin{tikzpicture}
        \node[
          fill=forest!10!paper, fill opacity=0.45,
          draw=forest, line width=0.9pt,
          text opacity=1, text=ink,
          rounded corners=5pt, font=\large,
          inner xsep=0.42cm, inner ysep=0.32cm,
          text width=0.90\linewidth, align=left
        ] {Can a few \textbf{collective variables} describe what depth does
           to information?};
      \end{tikzpicture}
    \end{column}

    \begin{column}{0.371\textwidth}
      \centering
      \ifanim
        \animategraphics[method=icon,autoplay,loop,poster=first,
          width=1.00\linewidth,type=png]%
          {20}{assets/elephant_frames/elephant_}{0}{59}%
      \else
        \includegraphics[width=1.00\linewidth]%
          {assets/von_neumann_elephant_wiggle_trail.png}%
      \fi

      \vspace{0.20em}
      \srcnote{Mayer, Khairy \& Howard,\\
      Am.\ J.\ Phys.\ \textbf{78}, 648 (2010)}
    \end{column}
  \end{columns}
  \endgroup
  \note{40 s. This is the physics thesis of the whole talk: trade millions
  of microscopic parameters for a handful of macroscopic ones. It sets up
  slide 7, where the reduction actually happens.}
\end{frame}

% =========================================================================
% 5. Broken telephone / information preservation  (~55 s)
% =========================================================================
\begin{frame}{At initialization, depth is like broken telephone}
  \begingroup
  \small
  \begin{center}
    Each layer \emph{re-transmits} the message it receives. Track the
    correlation between the two inputs, \(c_\ell\).

    \vspace{0.30em}
    {\color{forest}\bfseries Aim of the game: postpone perfect correlation
    for as long as possible!}
  \end{center}

  \vspace{0.80em}
  \begin{center}
  \makebox[\textwidth][c]{%
  \begin{tikzpicture}[x=1cm,y=1cm,>=Latex,scale=1.125,
                      every node/.style={transform shape}]
    % Signal paths first, so the layer bars sit on top of them.
    \foreach \y/\col in {0.95/orderedblue,-0.95/clay} {
      \draw[\col,line width=1.35pt,-{Latex[length=2.1mm]}]
        (1.58,\y)--(4.22,\y);
      \draw[\col,line width=1.35pt,-{Latex[length=2.1mm]}]
        (5.78,\y)--(8.42,\y);
    }

    % layer bars, two groups
    \foreach \x in {2.45,3.30,6.65,7.50} {
      \node[draw=forest!60,fill=forest!8,minimum width=0.32cm,
            minimum height=3.05cm,rounded corners=1pt] at (\x,0) {};
    }

    % --- column 1: the clean inputs ---
    \node[draw=orderedblue,line width=0.9pt,rounded corners=2pt,
          inner sep=1pt,fill=paper]
          at (0.90,0.95) {\telephoto{\catimg}{1.35cm}{orderedblue}{cat}};
    \node[draw=clay,line width=0.9pt,rounded corners=2pt,
          inner sep=1pt,fill=paper]
          at (0.90,-0.95) {\telephoto{\dogimg}{1.35cm}{clay}{dog}};
    \node[font=\bfseries\scriptsize,text=orderedblue,anchor=east]
          at (0.10,0.95) {cat \(\mathbf x^{(a)}\)};
    \node[font=\bfseries\scriptsize,text=clay,anchor=east]
          at (0.10,-0.95) {dog \(\mathbf x^{(b)}\)};

    % --- column 2: mid-depth, partly forgotten ---
    \node[draw=orderedblue!60,line width=0.9pt,rounded corners=2pt,
          inner sep=1pt,fill=paper]
          at (5.00,0.95) {\noisyphoto{\catimg}{1.35cm}{0.55}};
    \node[draw=clay!60,line width=0.9pt,rounded corners=2pt,
          inner sep=1pt,fill=paper]
          at (5.00,-0.95) {\noisyphoto{\dogimg}{1.35cm}{0.55}};

    % --- column 3: the SAME static both times. That is the point. ---
    \node[draw=muted,line width=0.9pt,rounded corners=2pt,
          inner sep=1pt,fill=paper]
          at (9.10,0.95) {\squarephoto{\staticimg}{1.35cm}{noise}};
    \node[draw=muted,line width=0.9pt,rounded corners=2pt,
          inner sep=1pt,fill=paper]
          at (9.10,-0.95) {\squarephoto{\staticimg}{1.35cm}{noise}};
    \node[font=\scriptsize,text=orderedblue,anchor=west]
          at (9.90,0.95) {\(\mathbf z_a^\ell\)};
    \node[font=\scriptsize,text=clay,anchor=west]
          at (9.90,-0.95) {\(\mathbf z_b^\ell\)};
    \draw[<->,muted,line width=0.7pt] (10.55,0.95)--(10.55,-0.95);
    \node[anchor=west,font=\scriptsize,text=muted,align=left]
          at (10.68,0) {same static:\\\(c_\ell\to1\)};

    % --- column captions ---
    \foreach \x/\lab in {0.90/{input},5.00/{mid-depth},9.10/{deep}} {
      \node[font=\scriptsize\color{muted}] at (\x,-1.95) {\lab};
    }
  \end{tikzpicture}}
  \end{center}

  \endgroup
  \note{55 s. This is the answer to the elephant slide: here is the
  collective variable. If \(c_\ell\to1\) the network has forgotten which
  input it saw, so deep layers cannot tell cats from dogs.}
\end{frame}

% =========================================================================
% 6. Three phases  (~50 s, 4 clicks)
% =========================================================================
\begin{frame}{Initialization sets how long distinctions survive}
  \begingroup
  \small
  \begin{columns}[T,onlytextwidth]
    \begin{column}{0.31\textwidth}
      \centering
      {\color{orderedblue}\bfseries Ordered}\\[0.10em]
      {\footnotesize \(c_\ell\to1\) fast;\\ everything looks alike}
    \end{column}
    \begin{column}{0.31\textwidth}
      \centering
      {\color{forest}\bfseries Critical edge}\\[0.10em]
      {\footnotesize forgetting is slow;\\ distinctions survive deep}
    \end{column}
    \begin{column}{0.31\textwidth}
      \centering
      {\color{clay}\bfseries Chaotic}\\[0.10em]
      {\footnotesize differences blow up,\\ then saturate}
    \end{column}
  \end{columns}

  \vspace{0.25em}
  \begin{center}
    \ifanim
      \animategraphics[method=icon,autoplay,loop,poster=first,
        width=0.77\textwidth,type=png]%
        {16}{assets/manifold_slide_frames_t/manifold_}{0}{200}%
    \else
      \only<1>{\includegraphics[width=0.77\textwidth]%
        {assets/manifold_slide_frames_t/manifold_0.png}}%
      \only<2>{\includegraphics[width=0.77\textwidth]%
        {assets/manifold_slide_frames_t/manifold_20.png}}%
      \only<3>{\includegraphics[width=0.77\textwidth]%
        {assets/manifold_slide_frames_t/manifold_80.png}}%
      \only<4->{\includegraphics[width=0.77\textwidth]%
        {assets/manifold_slide_frames_t/manifold_200.png}}%
    \fi

    \vspace{-0.10em}
    {\footnotesize\color{muted}A circular input manifold through random
    \(\tanh\) layers, projected to 2D.}

    \vspace{0.35em}
    \onslide<\manifoldbox->{%
    \begin{tikzpicture}
      \node[draw=forest,fill=forest!8,rounded corners=5pt,
            inner xsep=0.42cm,inner ysep=0.20cm,text=ink]
        {Tune initialization to the edge so distinctions propagate deepest.};
    \end{tikzpicture}}
  \end{center}
  \endgroup
  \note{Click through the four depths. The ordered panel collapsing to a
  point is the moment. 50 s.  Poole et al.\ 2016; Schoenholz et al.\ 2017.}
\end{frame}

% =========================================================================
% 7. Why the edge? Critical brain + the smoking gun  (~45 s)
% =========================================================================
\begin{frame}{The edge of chaos is a borrowed hypothesis}
  \begingroup
  \small
  \begin{columns}[T,onlytextwidth]
    \begin{column}{0.50\textwidth}
      {\color{forest}\bfseries The critical brain hypothesis}

      \vspace{0.10em}
      Cortical activity arrives in avalanches whose sizes follow a power
      law, \(P(s)\sim s^{-3/2}\). Scale-free all the way up to the
      recording array --- electrodes 200\,\(\mu\)m apart, millimetres of
      cortex firing as one event. No intrinsic scale, so the correlation
      length \(\xi\to\infty\).

      \vspace{0.15em}
      \srcnote{Beggs \& Plenz, J.\ Neurosci.\ \textbf{23}, 11167 (2003).
      Still contested: Beggs \& Timme, Front.\ Physiol.\ \textbf{3}, 163
      (2012).}

      \vspace{0.55em}
      {\color{forest}\bfseries Inherited by machine learning}

      \vspace{0.10em}
      {\footnotesize
      tuning nets to edge-of-chaos at initialization}

      \vspace{0.15em}
      \srcnote{Sompolinsky et al.\ 1988; Packard 1988;
      Bertschinger \& Natschl\"ager 2004; Poole et al.\ 2016;
      Schoenholz et al.\ 2017.}
    \end{column}

    \begin{column}{0.47\textwidth}
      \centering
      {\color{forest}\bfseries The smoking gun: \(\xi\) diverges}

      \vspace{0.20em}
      \begin{tikzpicture}
      \begin{axis}[
        width=5.0cm, height=2.85cm, scale only axis,
        xmin=0.55,xmax=1.45,ymin=0,ymax=27,
        xtick={0.6,0.8,1.0,1.2,1.4}, ytick={0,10,20},
        tick label style={font=\scriptsize,color=muted},
        axis line style={muted,line width=0.5pt},
        label style={font=\scriptsize,color=ink},
        xlabel={\(\chi\)}, ylabel={\(\xi_c\) (layers)},
        every axis plot/.append style={line width=1.0pt},
      ]
        \addplot[forest,dashed,line width=0.7pt,forget plot]
          coordinates {(1,0)(1,27)};
        \addplot[orderedblue,domain=0.55:0.962,samples=120]{1/(1-x)};
        \addplot[clay,domain=1.038:1.45,samples=120]{1/(x-1)};
        \node[font=\scriptsize,text=orderedblue] at (axis cs:0.72,6)
          {ordered};
        \node[font=\scriptsize,text=clay] at (axis cs:1.29,6) {chaotic};
        \node[font=\scriptsize\bfseries,text=forest,anchor=north]
          at (axis cs:1.0,26.5) {edge};
      \end{axis}
      \end{tikzpicture}

      \vspace{0.25em}
      {\footnotesize Criticality's signature is a \textbf{diverging correlation length} \(\xi\to\infty\)
      and power law behaviour.}

      \vspace{0.15em}
      {\footnotesize\color{muted}In a deep network that length is measured
      in \emph{layers}.}
    \end{column}
  \end{columns}

  \vspace{0.25em}
  \begin{center}
    \begin{tikzpicture}
      \node[draw=forest,fill=forest!8,rounded corners=5pt,
            inner xsep=0.38cm,inner ysep=0.16cm,text=ink]
        {Anything that cuts off the divergence destroys criticality
         \textbf{...enter dropout}};
    \end{tikzpicture}
  \end{center}
  \endgroup
  \note{45 s. Physics lineage, then the one diagnostic this room trusts.
  Say the critical-brain claim is contested -- do not oversell it. If
  pressed: Beggs \& Plenz 2003 is \emph{in vitro} (rat cortex cultures and
  acute slices, 8x8 array at 200 um pitch); in vivo support came later
  (Petermann et al., PNAS 2009). The -3/2 exponent is also just the
  mean-field critical-branching value, which is part of why the claim is
  contested. The box is the hinge into the second half of the talk.}
\end{frame}

% =========================================================================
% 8. The correlation length, shown  (~65 s)
% =========================================================================
\begin{frame}{Distance to perfect correlation decays exponentially}
  \begingroup
  \small
  \begin{center}
    How far are we from perfect correlation at layer \(\ell\)? Track
    \(|1-c_\ell|\): it decays exponentially, at a scale mean-field theory
    hands us explicitly:
    \vspace{-0.15em}
    \[
      |1-c_\ell|\sim e^{-\ell/\xi_c},
      \qquad
      \xi_c^{-1}=-\log\!\left[\sigma_w^2\!\int\!\!Dz_1Dz_2\,
        \phi'(u_1^*)\,\phi'(u_2^*)\right] .
    \]
    \vspace{-1.20em}

    {\footnotesize\color{muted}\(\sigma_w\) the weight std, \(\phi\) the
    nonlinearity, \(u_{1,2}^*\) the two inputs' preactivations at the fixed
    point. So \(\xi_c\) is set by \(\sigma_w,\sigma_b,\phi\) alone ---
    knobs we tune to make it \textbf{diverge}.}
  \end{center}

  \vspace{0.10em}
  \begin{center}
  \makebox[\textwidth][c]{%
  \begin{tikzpicture}[x=1cm,y=1cm,scale=0.87,
                      every node/.style={transform shape}]
    \foreach \x/\op/\cv in {0/0.00/0.10, 2.15/0.28/0.30, 4.30/0.48/0.50,
                            6.45/0.66/0.70, 8.60/0.83/0.90, 10.75/0.95/0.99} {
      \node[draw=orderedblue!70,line width=0.8pt,rounded corners=2pt,
            inner sep=0.8pt,fill=paper] at (\x,0.76)
            {\noisyphoto{\catimg}{1.24cm}{\op}};
      \node[draw=clay!70,line width=0.8pt,rounded corners=2pt,
            inner sep=0.8pt,fill=paper] at (\x,-0.76)
            {\noisyphoto{\dogimg}{1.24cm}{\op}};
      \node[font=\scriptsize,text=ink] at (\x,-1.68) {\(c_\ell=\cv\)};
    }
    \draw[muted,line width=0.7pt,-{Latex[length=2.0mm]}]
      (-0.70,-2.05)--(11.45,-2.05);
  \end{tikzpicture}}
  \end{center}

  \vspace{0.02em}
  \begin{center}
    {\color{forest}\bfseries Tune to the edge and \(\xi_c\to\infty\): the
    two never collapse together.}
  \end{center}
  \endgroup
  \note{65 s. The whole point in one picture: c grows toward 1 and the two
  inputs become the same thing. xi_c is how many layers that takes, and it
  is set by knobs we control. Full recursions and the propagator flow are
  in backup.}
\end{frame}

% =========================================================================
% 9. Dropout is an external field  (~75 s)
% =========================================================================
\begin{frame}{Dropout is an external field}
  \begingroup
  \small
  \begin{center}
  \makebox[\textwidth][c]{%
  \begin{tikzpicture}[x=1cm,y=1cm,scale=0.86,
                      every node/.style={transform shape}]
    % Same sweep as the previous slide, but each row now carries its OWN
    % dropout mask. The patterns never agree, so c saturates below 1.
    \foreach \x/\op/\cv in {0/0.00/0.10, 2.00/0.26/0.30, 4.00/0.46/0.50,
                            6.00/0.62/0.68, 8.00/0.74/0.80, 10.00/0.82/0.85} {
      \node[draw=orderedblue!70,line width=0.8pt,rounded corners=2pt,
            inner sep=0.8pt,fill=paper] at (\x,0.72)
        {\maskedphoto{\catimg}{1.16cm}{\op}%
          {0.07/0.62/0.24/0.19, 0.45/0.11/0.19/0.23, 0.71/0.63/0.21/0.22}};
      \node[draw=clay!70,line width=0.8pt,rounded corners=2pt,
            inner sep=0.8pt,fill=paper] at (\x,-0.72)
        {\maskedphoto{\dogimg}{1.16cm}{\op}%
          {0.58/0.56/0.26/0.22, 0.09/0.15/0.22/0.20, 0.35/0.73/0.19/0.18}};
      \node[font=\scriptsize,text=ink] at (\x,-1.58) {\(c_\ell=\cv\)};
    }
    \node[font=\scriptsize\bfseries,text=clay,anchor=west]
      at (11.05,0) {\(c^*=0.85<1\)};
    \draw[clay,line width=0.8pt] (10.95,-1.15) -- (10.95,1.15);
  \end{tikzpicture}}
  \end{center}

  \vspace{0.10em}
  \begin{center}
    {\footnotesize\color{muted}Each row gets its own mask, so even in pure
    noise the two never agree: \(c_\ell\) stalls below 1.}
  \end{center}

  \vspace{0.10em}
  \begin{center}
  \begin{minipage}{0.94\textwidth}
    {\color{clay}\bfseries \(c=1\) stops being a fixed point}, and
    \(\xi_c\) was a linearization \emph{about} it --- so it is not even
    defined. Look for the fixed point \emph{perturbatively away} from
    perfect correlation, in the dropout field \(h\propto1-\rho\).

    \onslide<2->{%
    \vspace{0.05em}
    Tuned to the edge (\(t\equiv\chi_\rho-1=0\)), one can show a length
    comes back:
    \vspace{-0.45em}
    {\Large\[
      \xi_c\ \propto\ h^{-1/2} .
    \]}}
  \end{minipage}
  \end{center}
  \endgroup
  \note{75 s. The dictionary is a 15 s gesture for this room; spend the
  time on the exponent table, which is the actual claim. Expect
  ``isn't this just mean field'' -- backup B3.}
\end{frame}

% =========================================================================
% 10. Two universality classes  (~40 s)
% =========================================================================
\begin{frame}{Two universality classes, neither of them Ising}
  \begingroup
  \small
  \begin{center}
    The critical exponents are set by the \textbf{analyticity of the
    activation}, not by a spatial dimension.
  \end{center}

  \vspace{0.60em}
  \begin{center}
    \renewcommand{\arraystretch}{1.35}
    \setlength{\tabcolsep}{10pt}
    \begin{tabular}{@{}lcc@{}}
      \toprule
      & \(\beta\) & \(\delta\) \\
      \midrule
      Smooth \(\phi\) (\(\tanh\)) & \(1\) & \(2\) \\
      Kinked \(\phi\) (\(\relu\)) & \(2\) & \(\tfrac32\) \\
      \midrule
      Ising, Landau mean field & \(\tfrac12\) & \(3\) \\
      \bottomrule
    \end{tabular}
  \end{center}

  \vspace{0.60em}
  \begin{center}
    \begin{tikzpicture}
      \node[draw=forest,fill=forest!8,rounded corners=5pt,
            inner xsep=0.40cm,inner ysep=0.18cm,text=ink]
        {A kink in \(\phi\) puts a branch point in \(\bar F_\rho\) at
         \(c=1\): \(m^{5/2}\) replaces \(m^{3}\).};
    \end{tikzpicture}
  \end{center}

  \vspace{0.45em}
  \scope{Structural analogy only: depth is an iteration count, not a
  distance, and \(m\ge0\) carries no \(\mathbb Z_2\) symmetry.}
  \endgroup
  \note{40 s. This is the paper's title claim. Full exponent table in
  backup; expect ``isn't this just mean field''.}
\end{frame}

% =========================================================================
% 10b. The ink-in-a-river intuition  (~40 s)
% =========================================================================
% One scene macro, driven by a single time parameter t in [0,1]: the ink
% plume advances downstream at the same current speed in all three rivers.
% \ifanim -> autoplaying animateinline; static build -> three overlay stills.
%   #1 = t
\newcommand{\riverscene}[1]{%
  \begin{tikzpicture}[x=1cm,y=1cm,>=Latex,scale=0.815,
                      every node/.style={transform shape}]
    % river = the signal flowing through depth; ink = dropout.
    % Same litre in all three rows, released in different places.
    \colorlet{rivercol}{orderedblue!16!paper}
    \colorlet{inkcol}{ink!84!orderedblue}
    \colorlet{bankcol}{orderedblue!40}
    % Frozen bounding box: every frame must occupy the same space or the
    % animation jitters as the plumes grow.
    \useasboundingbox (-2.15,-2.78) rectangle (11.05,3.34);

    \pgfmathsetmacro{\tt}{min(1,max(0,#1))}
    \pgfmathsetmacro{\vflow}{10.20}          % same current in every river
    \pgfmathsetmacro{\ftop}{min(11.60,8.40+\vflow*\tt)}
    \pgfmathsetmacro{\fmid}{min(11.60,0.80+\vflow*\tt)}
    \pgfmathsetmacro{\fbot}{min(11.60,1.40+\vflow*\tt)}
    % uniform row: concentration accumulates with distance, so the tint at
    % the plume tip deepens as the ramp fills out.
    \pgfmathsetmacro{\pct}{100*min(1,(\fmid-0.80)/9.90)}
    \colorlet{tipcol}{inkcol!\pct!rivercol}
    % the verdict labels fade in only once the plumes have run their course
    \pgfmathsetmacro{\labop}{max(0,min(1,(\tt-0.72)*4))}

    % ---- row 1: poured at the mouth -------------------------------------
    % banks are gentle sine curves so the bands read as water, not as bars.
    \pgfmathsetmacro{\btop}{min(9.20,\ftop)}
    \pgfmathsetmacro{\ctop}{max(\btop,\ftop-0.80)}
    \begin{scope}
      \clip plot[domain=0:10.70,samples=140,smooth]
              (\x,{2.34+0.05*sin(\x*105)})
            -- plot[domain=10.70:0,samples=140,smooth]
              (\x,{1.62+0.05*sin(\x*95+50)}) -- cycle;
      \fill[rivercol] (0,1.45) rectangle (10.70,2.50);
      \shade[left color=rivercol,right color=inkcol]
        (8.40,1.45) rectangle (\btop,2.50);
      \fill[inkcol] (\btop,1.45) rectangle (\ctop,2.50);
      \shade[left color=inkcol,right color=rivercol]
        (\ctop,1.45) rectangle (\ftop,2.50);
    \end{scope}
    \draw[bankcol,line width=0.6pt] plot[domain=0:10.70,samples=140,smooth]
      (\x,{2.34+0.05*sin(\x*105)});
    \draw[bankcol,line width=0.6pt] plot[domain=0:10.70,samples=140,smooth]
      (\x,{1.62+0.05*sin(\x*95+50)});
    \node[font=\scriptsize\bfseries,text=paper,text opacity=\labop]
      at (9.85,1.98) {last 20\%};

    % ---- row 2: poured evenly (the constant schedule) -------------------
    \pgfmathsetmacro{\emid}{min(11.20,\fmid+0.50)}
    \begin{scope}
      \clip plot[domain=0:10.70,samples=140,smooth]
              (\x,{0.62+0.05*sin(\x*98+20)})
            -- plot[domain=10.70:0,samples=140,smooth]
              (\x,{-0.10+0.05*sin(\x*112+65)}) -- cycle;
      \fill[rivercol] (0,-0.27) rectangle (10.70,0.78);
      % ink keeps being added along the way, so the tint ramps up downstream
      \shade[left color=rivercol,right color=tipcol]
        (0.80,-0.27) rectangle (\fmid,0.78);
      \shade[left color=tipcol,right color=rivercol]
        (\fmid,-0.27) rectangle (\emid,0.78);
    \end{scope}
    \draw[bankcol,line width=0.6pt] plot[domain=0:10.70,samples=140,smooth]
      (\x,{0.62+0.05*sin(\x*98+20)});
    \draw[bankcol,line width=0.6pt] plot[domain=0:10.70,samples=140,smooth]
      (\x,{-0.10+0.05*sin(\x*112+65)});
    \node[font=\scriptsize\bfseries,text=paper,text opacity=\labop]
      at (8.90,0.26) {all of it, but thin};

    % ---- row 3: poured upstream -----------------------------------------
    \pgfmathsetmacro{\bbot}{min(2.20,\fbot)}
    \pgfmathsetmacro{\cbot}{max(\bbot,\fbot-0.80)}
    \begin{scope}
      \clip plot[domain=0:10.70,samples=140,smooth]
              (\x,{-1.10+0.05*sin(\x*103+35)})
            -- plot[domain=10.70:0,samples=140,smooth]
              (\x,{-1.82+0.05*sin(\x*108+70)}) -- cycle;
      \fill[rivercol] (0,-1.99) rectangle (10.70,-0.94);
      \shade[left color=rivercol,right color=inkcol]
        (1.40,-1.99) rectangle (\bbot,-0.94);
      \fill[inkcol] (\bbot,-1.99) rectangle (\cbot,-0.94);
      \shade[left color=inkcol,right color=rivercol]
        (\cbot,-1.99) rectangle (\fbot,-0.94);
    \end{scope}
    \draw[bankcol,line width=0.6pt] plot[domain=0:10.70,samples=140,smooth]
      (\x,{-1.10+0.05*sin(\x*103+35)});
    \draw[bankcol,line width=0.6pt] plot[domain=0:10.70,samples=140,smooth]
      (\x,{-1.82+0.05*sin(\x*108+70)});
    \node[font=\scriptsize\bfseries,text=paper,text opacity=\labop]
      at (6.30,-1.46) {90\% of it, fully dark};

    % ---- the same litre, released three different ways ------------------
    \foreach \dx/\dy in {8.40/2.96, 1.40/-0.48} {
      \begin{scope}[shift={(\dx,\dy)}]
        \fill[inkcol] (0,0.27) .. controls (0.21,0.03) and (0.20,-0.12) ..
                      (0,-0.20) .. controls (-0.20,-0.12) and (-0.21,0.03) ..
                      cycle;
        \node[font=\tiny\bfseries,text=inkcol,anchor=west] at (0.30,0.02)
          {1\,L};
        \draw[inkcol,line width=0.9pt,-{Latex[length=1.7mm]}]
          (0,-0.32) -- (0,-0.56);
      \end{scope}
    }
    % the same litre, split into eight and sprinkled the whole way down
    \foreach \dx in {0.80,2.20,3.60,5.00,6.40,7.80,9.20,10.40} {
      \begin{scope}[shift={(\dx,1.24)},scale=0.60]
        \fill[inkcol] (0,0.27) .. controls (0.21,0.03) and (0.20,-0.12) ..
                      (0,-0.20) .. controls (-0.20,-0.12) and (-0.21,0.03) ..
                      cycle;
        \draw[inkcol,line width=0.9pt,-{Latex[length=1.7mm]}]
          (0,-0.34) -- (0,-0.72);
      \end{scope}
    }

    % ---- row labels ------------------------------------------------------
    \node[font=\scriptsize\bfseries,text=muted,anchor=east,align=right]
      at (-0.22,1.98) {poured at\\the mouth};
    \node[font=\scriptsize\bfseries,text=gilt,anchor=east,align=right]
      at (-0.22,0.26) {poured evenly\\(constant)};
    \node[font=\scriptsize\bfseries,text=forest,anchor=east,align=right]
      at (-0.22,-1.46) {poured\\upstream};

    % ---- flow axis -------------------------------------------------------
    \draw[muted,line width=0.7pt,-{Latex[length=2.0mm]}]
      (0,-2.24) -- (10.70,-2.24);
    \node[font=\scriptsize,text=muted,anchor=north west] at (0,-2.34)
      {input};
    \node[font=\scriptsize,text=muted,anchor=north east] at (10.70,-2.34)
      {output};
    \node[font=\scriptsize,text=muted,anchor=north] at (5.35,-2.34)
      {downstream: deeper layers \(\ell\)};
  \end{tikzpicture}%
}

\begin{frame}{Intuition: where would you pour the ink?}
  \begingroup
  \small
  \begin{center}
    Same litre of ink, same river, three ways to pour it in.
  \end{center}

  \vspace{0.05em}
  \begin{center}
  \makebox[\textwidth][c]{%
  \ifanim
    % 30 frames at 15 fps = 2 s; t saturates at frame 24 so the finished
    % picture is held for the last third of the loop.
    \begin{animateinline}[autoplay,loop,poster=last]{15}
      \multiframe{30}{i=0+1}{\riverscene{\i/24}}
    \end{animateinline}%
  \else
    \only<1>{\riverscene{0.06}}%
    \only<2>{\riverscene{0.34}}%
    \only<3->{\riverscene{1}}%
  \fi
  }
  \end{center}

  \vspace{0.05em}
  \begin{center}
    {\color{forest}\bfseries Dropout is the ink, depth is the river.}
    Constant dropout is the middle one --- front-loading is the bottom one.
  \end{center}
  \endgroup
  \note{40 s. Pure intuition slide, no equations. The budget constraint is
  the litre: you do not get to add ink, only to choose where it goes.
  Everything downstream of the release point carries the perturbation, so
  early placement buys the most reach per unit of dropout, and spreading it
  evenly wastes half the river on a tint too thin to matter. The next slide
  is this argument done properly, as a variational problem.}
\end{frame}

% =========================================================================
% 11. Optimal schedule  (~75 s)
% =========================================================================
\begin{frame}{Optimal schedule: saturated blocks, placed early}
  \begingroup
  \small
  \begin{columns}[T,onlytextwidth]
    \begin{column}{0.52\textwidth}
      Depth acts like an RG scale, so the couplings \emph{run}: the dropout
      field becomes a function of depth, \(h\to h_\ell\).

      \vspace{0.45em}
      Treat the schedule as a functional and maximize the correlation
      length at fixed budget:
      \vspace{-0.10em}
      \[
        \max_{\{h_\ell\}}\ \xi_{\mathrm{eff}}[h]
        \quad\text{s.t.}\quad
        \frac1L\sum_\ell h_\ell=\bar h,\ \ 0\le h_\ell\le h_{\max}.
      \]

      \vspace{-0.10em}
      \onslide<2->{%
      At the edge \(\xi_{\mathrm{eff}}^{-1}[h]\propto
      \frac1L\sum_\ell h_\ell^{\,a}\) with \(a<1\), so maximizing
      \(\xi_{\mathrm{eff}}\) means minimizing a \textbf{concave} sum at
      fixed mean. The optimum sits at the corners:
      \(h_\ell\in\{0,h_{\max}\}\).}

      \vspace{0.40em}
      \onslide<3->{%
      \begin{tikzpicture}
        \node[draw=forest,fill=forest!8,rounded corners=5pt,
              inner xsep=0.26cm,inner ysep=0.14cm,text=ink,
              text width=0.91\linewidth,font=\footnotesize]
          {\textbf{Propagation fixes the shape.} It is order-invariant, so
           it cannot say \emph{where} the block goes.\\[0.25em]
           \textbf{Reach fixes the position.} An early mask perturbs a
           longer downstream suffix.};
      \end{tikzpicture}}
    \end{column}

    \begin{column}{0.44\textwidth}
      \centering
      \includegraphics[width=0.86\linewidth]%
        {assets/results/frontloaded_schedule_comparison.pdf}

      \vspace{0.10em}
      {\footnotesize\color{muted}Three profiles, identical
      \(\sum_\ell h_\ell\).}
    \end{column}
  \end{columns}

  \vspace{0.05em}
  \onslide<3->{\scope{Reach is a separate argument, and the weaker of the two.}}
  \endgroup
  \note{75 s. Be explicit that this is a two-part argument -- do not paper
  over the seam. If pushed on reach, concede it is the softer half.}
\end{frame}

% =========================================================================
% 12. Experiments  (~75 s)
% =========================================================================
\begin{frame}{The predicted schedule wins, and wins sooner}
  \begingroup
  \small
  \begin{columns}[c,onlytextwidth]
    \begin{column}{0.54\textwidth}
      \centering
      {\color{forest}\bfseries Best test accuracy across schedules}\\[0.05em]
      {\footnotesize CIFAR-10, \(\relu\) MLP, width 256}\\[0.15em]
      \includegraphics[width=0.94\linewidth]%
        {assets/results/h_sweep_schedule_styled.pdf}

      \vspace{0.10em}
      {\footnotesize\color{muted}Mean \(\pm\) s.e.\ over \mlpruns\ runs.}
    \end{column}

    \begin{column}{0.42\textwidth}
      \centering
      {\color{forest}\bfseries Final-epoch test \emph{loss}}\\[0.05em]
      {\fontsize{19}{22}\selectfont\color{clay}\bfseries 18--35\% lower}\\
      {\footnotesize MLP settings}

      \vspace{0.22em}
      {\color{gold}\rule{4.2cm}{1.0pt}}

      \vspace{0.22em}
      {\fontsize{19}{22}\selectfont\color{clay}\bfseries 4--6\% lower}\\
      {\footnotesize Vision Transformers}

      \vspace{0.22em}
      {\color{gold}\rule{4.2cm}{1.0pt}}

      \vspace{0.22em}
      {\color{forest}\bfseries Epochs to match constant dropout}\\[0.05em]
      {\fontsize{19}{22}\selectfont\color{clay}\bfseries
       \epochsfast\ vs \epochsbase}\\
      {\footnotesize \(\approx\)60\% fewer, at equal per-epoch cost}
    \end{column}
  \end{columns}

  \vspace{0.25em}
  \scope{Left panel is accuracy; the percentages are cross-entropy.
  ``Fewer epochs'' is epoch-proportional compute at identical architecture
  and batch size, not wall-clock or per-step FLOPs.}
  \endgroup
  \note{75 s. Say ``fewer epochs'', never ``less compute'' -- the honest
  claim is just as good. Loss vs accuracy distinction out loud.
  Confirm \mlpruns\ runs against the archived config before presenting.}
\end{frame}

% =========================================================================
% 13. Close  (~30 s)
% =========================================================================
\begin{frame}[plain]
  \begin{tikzpicture}[remember picture,overlay]
    \fill[paper] (current page.south west) rectangle (current page.north east);
    \node[anchor=north west] at ($(current page.north west)+(0.75,-0.38)$)
      {\includegraphics[height=0.82cm]{\iaifilogo}};
    \node[anchor=north east] at ($(current page.north east)+(-0.30,-0.18)$)
      {\includegraphics[height=1.74cm]{\cmulogo}};
  \end{tikzpicture}

  \vspace{1.15cm}
  \begin{center}
    \resizebox{0.70\paperwidth}{!}{%
      \color{forest}\bfseries From microscopic masks to a macroscopic design rule}

    \vspace{0.52cm}
    {\Large Dropout hurts propagation.}\\[0.25em]
    {\color{gold}\Large\(\Downarrow\)}\\[0.20em]
    {\Large Maximize the correlation length.}\\[0.25em]
    {\color{gold}\Large\(\Downarrow\)}\\[0.20em]
    {\Large\color{forest}\bfseries Front-load dropout.}

    \vspace{0.45cm}
    {\color{gold}\rule{6.5cm}{1.0pt}}

    \vspace{0.35cm}
    {\small\color{muted}Outlook: treat other depth-dependent
    hyperparameters as dynamical fields.}

    \vspace{0.32cm}
    {\large\color{forest}\bfseries Questions?}\\[0.15em]
    {\footnotesize\color{muted}arXiv:2605.21648 \(\cdot\)
     Carnegie Mellon University}
  \end{center}
\end{frame}

% #########################################################################
% BACKUP
% #########################################################################

% ---- B-2: mean-field propagator flow (moved off the main path) ---------
\begin{frame}{Backup: wide networks collapse to one map in depth}
  \begingroup
  \small
  \begin{columns}[T,onlytextwidth]
    \begin{column}{0.49\textwidth}
      At \(N\to\infty\) the layer-\(\ell\) theory is \emph{free}: the action
      is quadratic, so the propagator \(q^\ell_{ab}\) is its only datum.
      \vspace{-0.10em}
      \[
        p(z^\ell)\propto e^{-S_\ell[z]},\quad
        S_\ell=\tfrac12\textstyle\sum_{i,a,b}\big(q_\ell^{-1}\big)_{ab}
          z_{i;a}z_{i;b}
      \]

      \vspace{-0.25em}
      Integrating out a layer makes it \emph{flow}:
      \vspace{-0.15em}
      \[
        q^{\ell+1}_{ab}=\sigma_w^2\big\langle\phi(z_a)\phi(z_b)
          \big\rangle_{q_\ell}+\sigma_b^2 .
      \]

      \vspace{-0.20em}
      Once \(q^\ell_{aa}\to q^*\), only \(c_\ell=q^\ell_{ab}/q^*\) survives:
      \vspace{-0.15em}
      \[
        \begin{gathered}
          \boxed{\;c_{\ell+1}=F(c_\ell)\;}\\[0.20em]
          F(c)=\big[\sigma_w^2\big\langle\phi(u_a)\phi(u_b)\big\rangle_c
            +\sigma_b^2\big]\big/q^*
        \end{gathered}
      \]

      \vspace{-0.20em}
      {\footnotesize\color{muted}\((u_a,u_b)\) Gaussian,
      \(\E[u_au_b]=c\,q^*\).}
    \end{column}

    \begin{column}{0.48\textwidth}
      \onslide<2->{%
      \centering
      \begin{tikzpicture}
      \begin{axis}[
        width=5.0cm, height=2.10cm, scale only axis,
        xmin=0,xmax=40,ymin=0,ymax=1.14,
        xtick={0,10,20,30,40},ytick={0,0.5,1},
        tick label style={font=\scriptsize,color=muted},
        axis line style={muted,line width=0.5pt},
        label style={font=\scriptsize,color=ink},
        xlabel={depth \(\ell\)}, ylabel={overlap \(c_\ell\)},
        every axis plot/.append style={line width=1.0pt},
        legend style={at={(0.97,0.05)},anchor=south east,font=\scriptsize,
                      draw=none,fill=none,row sep=-1pt},
        legend cell align=left,
      ]
        % "forget plot" keeps this reference line out of the legend; without
        % it pgfplots shifts every entry down by one row.
        \addplot[muted,dashed,line width=0.6pt,domain=0:40,samples=2,
                 forget plot]{1};
        \addplot[orderedblue,domain=0:40,samples=160]{1-0.8*pow(0.5,x)};
          \addlegendentry{\(\lambda=0.5\)}
        \addplot[gilt,domain=0:40,samples=160]{1-0.8*pow(0.85,x)};
          \addlegendentry{\(\lambda=0.85\)}
        \addplot[forest,domain=0:40,samples=160]{1-0.8*pow(0.98,x)};
          \addlegendentry{\(\lambda=0.98\)}
        \node[font=\scriptsize,text=muted,anchor=south east]
          at (axis cs:40,1.01) {\(c^*=1\)};
      \end{axis}
      \end{tikzpicture}

      \vspace{0.20em}
      {\footnotesize\color{muted}Same fixed point \(c^*=F(c^*)\);
      only the rate differs.}

      \vspace{0.35em}
      The slope \(\lambda=F'(c^*)\) sets an e-folding depth
      \vspace{-0.20em}
      \[
        \xi_c=-1/\log|\lambda| \quad\text{layers.}
      \]
      \vspace{-1.05em}
      {\footnotesize\color{muted}Here \(1.4,\ 6,\ 50\).}}
    \end{column}
  \end{columns}

  \vspace{0.10em}
  \begin{center}
    {\color{forest}\bfseries Millions of weights become a one-dimensional
    dynamical system in depth.}
  \end{center}

  \vspace{0.05em}
  \scope{The quadratic action is exact only at \(N=\infty\); finite width
  generates interactions at \(\mathcal O(1/N)\). Without dropout
  \(\lambda=F'(1)\equiv\chi\).}
  \endgroup
  \note{75 s. The cobweb replaces the Schoenholz heatmap: same message,
  parsed instantly. Full Gaussian recursions are in backup.}
\end{frame}

% ---- B-1: double descent (moved off the main path) ---------------------
\begin{frame}{Backup: regularization lost its textbook justification}
  \begingroup
  \small
  \begin{center}
  \begin{tikzpicture}
  \begin{axis}[
    width=8.6cm, height=3.45cm, scale only axis,
    xmin=0, xmax=10.6, ymin=0, ymax=1.42,
    xtick=\empty, ytick=\empty,
    axis lines=left,
    axis line style={-{Latex[length=2.0mm]},muted,line width=0.8pt},
    xlabel={increasing model capacity},
    ylabel={error},
    label style={font=\footnotesize,color=ink},
    every axis x label/.style={at={(axis description cs:0.78,-0.06)},
                               anchor=north},
    every axis y label/.style={at={(axis description cs:-0.03,0.55)},
                               anchor=south,rotate=90},
    clip=false,
  ]
    % --- regime shading, colours matched to the two column headers below ---
    \fill[gilt!20!paper] (axis cs:0,0) rectangle (axis cs:5,1.42);
    \only<2->{\fill[forest!16!paper]
      (axis cs:5,0) rectangle (axis cs:10.6,1.42);}

    % --- test error ---
    \addplot[ink,line width=1.15pt,domain=0:5,samples=140]
      {0.62*exp(-1.3*x)+0.20+0.9*(x/5)^4};
    \only<2->{\addplot[ink,line width=1.15pt,domain=5:10.6,samples=140]
      {0.15+0.95*exp(-1.1*(x-5))};}

    % --- train error ---
    \addplot[muted,dashed,line width=1.0pt,domain=0:5,samples=120]
      {0.55*exp(-0.9*x)};
    \only<2->{\addplot[muted,dashed,line width=1.0pt,domain=5:10.6,
      samples=120]{0.55*exp(-0.9*x)};}

    % --- interpolation threshold ---
    \draw[muted,dashed,line width=0.7pt]
      (axis cs:5,0) -- (axis cs:5,1.14);
    \node[font=\scriptsize,text=muted,rotate=90,anchor=south]
      at (axis cs:4.90,0.30) {interpolation threshold};

    % --- curve labels ---
    \node[font=\scriptsize,text=ink,rotate=-40,anchor=south west]
      at (axis cs:0.55,0.60) {test error};
    \node[font=\scriptsize,text=muted,rotate=-32]
      at (axis cs:1.05,0.17) {train error};

    % --- regime captions ---
    \node[font=\scriptsize,text=gilt,align=center]
      at (axis cs:2.45,1.28) {classical regime\\(underparametrized)};
    \only<2->{\node[font=\scriptsize,text=forest,align=center]
      at (axis cs:7.9,1.28) {modern interpolating regime\\(overparametrized)};}
  \end{axis}
  \end{tikzpicture}
  \end{center}

  \vspace{0.25em}
  \begin{columns}[T,onlytextwidth]
    \begin{column}{0.47\textwidth}
      {\color{gilt}\bfseries Classical regime}

      \vspace{0.10em}
      More capacity overfits, so a regularizer earns its keep by
      \emph{removing} capacity.
    \end{column}

    \begin{column}{0.49\textwidth}
      \onslide<2->{%
      {\color{forest}\bfseries Modern interpolating regime}

      \vspace{0.10em}
      Past the threshold more capacity \emph{helps}. Our networks live
      here, where that story fails.}
    \end{column}
  \end{columns}

  \vspace{0.25em}
  \onslide<2->{%
  \begin{center}
    \begin{tikzpicture}
      \node[draw=forest,fill=forest!8,rounded corners=5pt,
            inner xsep=0.40cm,inner ysep=0.16cm,text=ink]
        {So what \emph{is} dropout doing, how much of it, and
         \textbf{where in depth}?};
    \end{tikzpicture}
  \end{center}}
  \endgroup
  \note{50 s. Opening move: the textbook reason for regularizing does not
  apply to the models we train. That is the gap the rest of the talk fills
  with criticality rather than capacity. Curve redrawn after
  Belkin et al., PNAS 116, 15849 (2019).}
\end{frame}

% ---- B0: scaling collapse (moved off the main path) --------------------
\begin{frame}{Backup: two parameters collapse onto one scaling variable}
  \begingroup
  \centering
  \vspace{-0.30em}
  \includegraphics[width=0.74\textwidth]%
    {../figures/theory/smooth_scaling_collapse.pdf}

  \vspace{0.20em}
  {\small
   smooth \(m=h^{1/2}\mathcal M_s\!\left(t/h^{1/2}\right)\)
   \qquad\(\cdot\)\qquad
   kinked \(m=h^{2/3}\mathcal M_k\!\left(t/h^{1/3}\right)\)}

  \vspace{0.20em}
  {\small Solid line is the parameter-free normal form
   \(\tilde m=\sqrt{1+\tilde t^{\,2}}-\tilde t\), not a fit.}

  \vspace{0.15em}
  \scope{Collapse holds in the scaling region; the spread at
  \(\tilde t\lesssim-0.5\) is corrections to scaling at the largest \(h\).}
  \endgroup
\end{frame}

% ---- B1: does this survive training? (most likely question) -------------
\begin{frame}{Backup: the theory is at initialization; the results are trained}
  \begingroup
  \small
  {\color{clay}\bfseries The gap, stated plainly.}
  \(\xi_{\mathrm{eff}}\) is a property of the \emph{forward map at
  initialization}. The experiments measure test accuracy after full
  training. Nothing here derives the second from the first.

  \vspace{0.55em}
  {\color{forest}\bfseries What is actually claimed}
  \begin{itemize}
    \setlength{\itemsep}{0.30em}
    \item The schedule shape is chosen \emph{a priori} from the
      initialization theory, with no tuning on the test set.
    \item That prediction, made in advance, is the one that wins.
    \item So the initialization functional is a useful \emph{design
      heuristic} for a schedule, not a theorem about the trained network.
  \end{itemize}

  \vspace{0.55em}
  {\color{forest}\bfseries Why it is not vacuous.} The ordering of
  schedules is predicted before any training run, and the predicted
  ordering is reproduced across MLPs and ViTs, two architectures with
  different \(F\).

  \vspace{0.45em}
  \scope{Honest answer if pressed: connecting init-time \(\xi_c\) to
  end-of-training generalization is open, and is the main thing I would
  want to do next.}
  \endgroup
\end{frame}

% ---- B2: isn't this just mean field? -----------------------------------
\begin{frame}{Backup: ``isn't this all just mean field?''}
  \begingroup
  \small
  {\color{forest}\bfseries Yes, and that is not the interesting axis.}

  \vspace{0.40em}
  \begin{itemize}
    \setlength{\itemsep}{0.35em}
    \item \(N\to\infty\) is exactly a mean-field limit. There are no
      fluctuation corrections at leading order.
    \item There is \textbf{no spatial dimension} here. Depth is an
      iteration index, not a coordinate, so there is no \(d\), no upper
      critical dimension, and no \(\epsilon\)-expansion.
    \item The universality is therefore not over dimension. It is over
      the \textbf{analyticity class of \(\phi\)} at the fixed point.
  \end{itemize}

  \vspace{0.55em}
  {\color{forest}\bfseries Why the exponents still differ from Landau.}
  A kink in \(\phi\) puts a branch point in \(\bar F_\rho\) at \(c=1\).
  The effective potential picks up \(m^{5/2}\) rather than \(m^{3}\), and
  the exponents follow from that non-analyticity. Same mean-field
  machinery, different local expansion.

  \vspace{0.45em}
  \scope{So ``mean field'' describes the method; it does not determine
  the class.}
  \endgroup
\end{frame}

% ---- B3: finite width --------------------------------------------------
\begin{frame}{Backup: finite width}
  \begingroup
  \small
  \begin{itemize}
    \setlength{\itemsep}{0.40em}
    \item Finite \(N\) makes \(c_\ell\) a random variable with
      \(\mathcal O(1/N)\) variance about the deterministic recursion.
    \item Near the edge, \(\lambda\to1\), so fluctuations accumulate over
      depth rather than damping. The width needed for the mean-field
      description to hold grows with \(L\).
    \item Experiments run at width 256 with \(L\) of order tens, where the
      recursion tracks the empirical overlap well.
    \item The scheduling conclusion depends only on \emph{concavity} of
      \(h\mapsto h^a\), which is robust to the fluctuation correction
      shifting \(a\) slightly.
  \end{itemize}

  \vspace{0.55em}
  \scope{The 1/N corrections are not computed in this work. Roberts,
  Yaida \& Hanin (2022) develop the systematic expansion.}
  \endgroup
\end{frame}

% ---- B4: why do ViTs transfer? -----------------------------------------
\begin{frame}{Backup: why the ViT numbers transfer}
  \begingroup
  \small
  {\color{clay}\bfseries The fair objection.} Attention, LayerNorm and
  residual connections all change \(F\). The MLP derivation does not
  literally apply.

  \vspace{0.50em}
  {\color{forest}\bfseries What survives the change of \(F\)}
  \begin{itemize}
    \setlength{\itemsep}{0.30em}
    \item The argument needs only two things: dropout acts as a
      decorrelating field \(h>0\), and \(\xi_{\mathrm{eff}}^{-1}\) is
      \emph{concave} in \(h\).
    \item Neither depends on the specific form of \(F\), only on the
      existence of a fixed point at \(c=1\) that dropout displaces.
    \item So the prediction is qualitative for ViTs: front-load, whatever
      the exponent \(a\) turns out to be.
  \end{itemize}

  \vspace{0.50em}
  {\color{forest}\bfseries Observed.} 4--6\% final-epoch loss reduction,
  same ordering of schedules. Component ablations put dropout on the
  attention block, the MLP block, and both; the early-step schedule wins
  in every case.

  \vspace{0.40em}
  \scope{The ViT exponent is not measured. Only the sign of the effect is
  predicted.}
  \endgroup
\end{frame}

% ---- B5: full mean-field recursions ------------------------------------
\begin{frame}{Backup: the full second-moment recursions}
  \begingroup
  \small
  \begin{columns}[T,onlytextwidth]
    \begin{column}{0.62\textwidth}
      \begin{align*}
        z_{i;a}^{\ell+1}
          &=\sum_{j=1}^{N}W_{ij}^{\ell+1}\phi(z_{j;a}^{\ell})+b_i^{\ell+1},\\
        W_{ij}^{\ell}&\sim\mathcal N(0,\sigma_w^2/N),\quad
        b_i^{\ell}\sim\mathcal N(0,\sigma_b^2).
      \end{align*}
      \begin{align*}
        q_{aa}^{\ell+1}
          &=\sigma_w^2\!\int\! Dz\,\phi^2\!\left(\sqrt{q_{aa}^{\ell}}z\right)
            +\sigma_b^2,\\
        q_{ab}^{\ell+1}
          &=\sigma_w^2\,\E_{(u_a,u_b)}[\phi(u_a)\phi(u_b)]+\sigma_b^2,\\
        c_{ab}^{\ell}
          &=q_{ab}^{\ell}\big/\sqrt{q_{aa}^{\ell}q_{bb}^{\ell}}.
      \end{align*}
    \end{column}
    \begin{column}{0.34\textwidth}
      {\color{forest}\bfseries Assumptions}
      \begin{itemize}
        \setlength{\itemsep}{0.30em}
        \item \(N\to\infty\)
        \item i.i.d.\ Gaussian init
        \item fixed activation \(\phi\)
        \item exchangeable neurons
      \end{itemize}

      \vspace{0.50em}
      {\footnotesize \((u_a,u_b)\) jointly Gaussian with covariance fixed
      by \(q_{aa}^\ell,q_{bb}^\ell,q_{ab}^\ell\); \(Dz\) is the standard
      Gaussian measure.}

      \vspace{0.40em}
      {\footnotesize At equal variance,
      \(\E\|\mathbf z_a-\mathbf z_b\|^2/N=2q^*(1-c)\).}
    \end{column}
  \end{columns}
  \endgroup
\end{frame}

% ---- B6: RG coefficient table (from old slide 10) ----------------------
\begin{frame}{Backup: depth as an RG scale}
  \begingroup
  \small
  \begin{center}
    {\bfseries network depth \(\ell\uparrow\)
      \(\longleftrightarrow\) QFT scale \(\mu\downarrow\)
      toward the infrared}
  \end{center}

  \vspace{0.35em}
  Coarse-graining representations makes the effective coefficients run
  with depth:

  \vspace{0.25em}
  \begin{center}
    \footnotesize
    \setlength{\tabcolsep}{6pt}
    \renewcommand{\arraystretch}{1.20}
    \begin{tabular}{@{}lll@{}}
      \toprule
      \textbf{coefficient} & \textbf{network definition} &
      \textbf{effective-theory role} \\
      \midrule
      \(t_\ell\) & \(\chi_{\rho_\ell}-1\) & mass-like detuning \\
      \(h_\ell\) & \(1-\bar F_{\rho_\ell}(1)\) & relevant decorrelation source \\
      \(g_\ell\) & \(\bar F_{\rho_\ell}''(1)\) & smooth analytic interaction \\
      \(\kappa_\ell\) & branch-point amplitude & kinked non-analytic interaction \\
      \bottomrule
    \end{tabular}
  \end{center}

  \vspace{0.40em}
  Both equations of state follow from a reduced (zero-momentum) 1PI action
  \[
    \Gamma_{\mathrm{eff}}=-\tfrac{t}{2}m^2+\mathcal V(m)-hm,
    \qquad
    \mathcal V_s=\tfrac{g_\rho}{6}m^3,
    \qquad
    \mathcal V_k=\tfrac{2\kappa}{5}m^{5/2}.
  \]

  \vspace{0.30em}
  \scope{RG-scale analogy. Generic MLPs are not thereby proven conformal,
  and \(\Gamma_{\mathrm{eff}}\) is not the training loss or a neural
  partition function. Roberts, Yaida \& Hanin (2022), Ch.~4.}
  \endgroup
\end{frame}

% ---- B7: kinked collapse -----------------------------------------------
\begin{frame}{Backup: the kinked (\(\relu\)) collapse}
  \begingroup
  \centering
  \includegraphics[width=0.82\textwidth]%
    {../figures/theory/kinked_scaling_collapse.pdf}

  \vspace{0.30em}
  \small
  \(m_k(t,h)=h^{2/3}\,\mathcal M_k(t/h^{1/3})\), with
  \(\mathcal V_k=\tfrac{2\kappa}{5}m^{5/2}\) from the branch point at
  \(c=1\).

  \vspace{0.25em}
  \scope{Ordinary bias-free \(\relu\) follows the constrained path
  \(t=-h\) and retains \(\xi\sim h^{-1/3}\).}
  \endgroup
\end{frame}

% ---- B8: full exponent dictionary --------------------------------------
\begin{frame}{Backup: the complete critical-exponent dictionary}
  \begin{center}
    \small
    \setlength{\tabcolsep}{3.5pt}
    \renewcommand{\arraystretch}{1.10}
    \begin{tabular}{@{}lccccccc@{}}
      \toprule
      Universality class & \(\nu_t\) & \(\beta\) & \(\theta_{\rm rel}\) &
      \(\gamma\) & \(\delta\) & \(\nu_\rho\) & \(\alpha\) \\
      \midrule
      Smooth activation & \(1\) & \(1\) & \(1\) & \(1\) & \(2\) & \(\tfrac12\) & \(-1\) \\
      Kinked activation & \(1\) & \(2\) & \(2\) & \(1\) & \(\tfrac32\) & \(\tfrac13\) & \(-3\) \\
      Spin glass (SK MF) & \(\tfrac12\) & \(1\) & \(2\) & \(1\) & \(2\) & \(\tfrac14\) & \(-1\) \\
      Ising-like (Landau MF) & \(\tfrac12\) & \(\tfrac12\) & \(1\) & \(1\) & \(3\) & \(\tfrac13\) & \(0\) \\
      \bottomrule
    \end{tabular}
  \end{center}

  \vspace{0.30em}
  \begin{columns}[T,onlytextwidth]
    \begin{column}{0.48\textwidth}
      \small
      \setlength{\jot}{1.5pt}
      \begin{align*}
        m&\sim |t|^\beta &&(h=0),\\
        m&\sim |h|^{1/\delta} &&(t=0),\\
        \chi_M&\sim |t|^{-\gamma},\\
        C&\sim |t|^{-\alpha}.
      \end{align*}
    \end{column}
    \begin{column}{0.48\textwidth}
      \small
      \setlength{\jot}{1.5pt}
      \begin{align*}
        \xi&\sim |t|^{-\nu_t} &&(h=0),\\
        \xi&\sim |h|^{-\nu_\rho} &&(t=0),\\
        m_\ell&\sim \ell^{-\theta_{\rm rel}} &&(t=h=0).
      \end{align*}
    \end{column}
  \end{columns}

  \vspace{0.20em}
  \scope{The kinked \(\nu_\rho\) is a normal-form field exponent.}
\end{frame}

% ---- B9: prior work ----------------------------------------------------
\begin{frame}{Backup: relation to prior scheduled-dropout work}
  \begingroup
  \small
  \begin{itemize}
    \setlength{\itemsep}{0.40em}
    \item \textbf{Stochastic depth} (Huang et al., 2016) drops whole
      residual blocks with a probability that increases with depth. A
      mechanism, not a derived profile.
    \item \textbf{Curriculum dropout} (Morerio et al., 2017) anneals
      dropout over \emph{training time}. Orthogonal axis: ours is a
      profile over \emph{layers}.
    \item \textbf{LayerDrop} (Fan et al., 2019) removes transformer layers
      for inference efficiency.
    \item \textbf{Hayou et al.\ (2019)} observed that smooth and
      \(\relu\)-like activations behave differently at the edge of chaos.
  \end{itemize}

  \vspace{0.55em}
  {\color{forest}\bfseries What is new here.} The exponents are extracted
  and the two universality classes identified, and the schedule shape is
  \emph{derived} from concavity of the correlation-depth functional rather
  than proposed and then tuned.
  \endgroup
\end{frame}

\end{document}
